\section{Discussion and limitations}
The representation supplies a direct relation between spherical geometry, angular resolution, and physical noise. Multipole coefficients expose reconstruction errors that total-state losses and global Husimi correlation can conceal. The controlled experiment also changes the interpretation of an initial missing-mode result: failure at a short training budget does not establish an expressivity limit of the channel family.

The scope remains deliberately small. The data consist of one synthetic mixture family. Single-spin errors are empirical reconstruction errors, and the confirmation split measures repeatability across newly sampled datasets rather than out-of-sample generalization of a frozen model. Finite-$j$ encoding discards high harmonics and attenuates lower ones; even exact density reconstruction generates the smoothed readout, not the original density. A grid-dependent mode criterion adds further resolution dependence.

Several comparisons are necessary before attributing a generative benefit to diffusion. Final-only training uses a composition of physical channels but no intermediate forward targets. A matched single-channel or direct state-preparation baseline, a normalized path loss, and a learned classical spherical diffusion baseline would distinguish the role of the forward construction from that of parameterized state preparation. The present data cannot demonstrate that the six-stage structure is necessary. Likewise, the forward spectrum orders decay but does not establish coarse-to-fine ordering of the learned reverse dynamics.

The correlated pilot is a clear current limitation. Additional reverse capacity, improved supervision, or different ancilla designs might help, but no such improvement is claimed without experiments. Multiple independent train/test splits and jointly reported state, correlation, parameter, compute and sampling costs are needed. Classical correlations in the source already admit a separable encoded representation, so an exponential Hilbert-space dimension alone is not evidence of a quantum computational advantage.

Hardware execution is also untested. The spin-$j$ space can be represented as the symmetric sector of $2j$ qubits, but a physical device requires specific preparation, control, reset, and readout implementations. The simulated coherent-state POVM and dense matrix trajectory do not supply gate-depth estimates or a scalable sampling procedure. The work supports an operationally quantum model implemented in simulation, not a demonstrated hardware-native generator.

\section{Conclusion}
Spin coherent-state encoding and isotropic Lindblad diffusion give a finite-dimensional generative construction with an exact spherical heat spectrum. The numerical evidence separates representation limits from training failures: final-only supervision restores the representable secondary mode without enlarging the reverse architecture, and validation-selected multipole weighting improves high-rank reconstruction on new seeds. The poor two-spin pilot and the absence of direct-preparation and learned classical diffusion comparisons define the immediate limits of the result. The accompanying evidence archive makes these positive and negative conclusions inspectable and provides a stable basis for extending the study.

\section*{Data and code availability}
The accompanying project includes frozen density matrices, parameters, configurations, study protocols, source snapshots, checksum manifests, and scripts that audit the principal results and regenerate every paper table and figure. A claim-to-evidence ledger identifies the support and limitations of each result. The 70-run diagnostic/hybrid source snapshot is contemporaneous; earlier studies have less complete source provenance, as recorded in the archive. No public repository URL or archival DOI has yet been designated. Author details, funding, contributions, competing-interest declarations, and a permanent public archive must be completed by the authors before submission.
